\documentclass[11pt,a4paper]{article}
\usepackage{booktabs}
\usepackage{caption}
\usepackage{multirow}
\usepackage{makecell}
\usepackage{amsmath}
\usepackage{amssymb}
\usepackage{array}
\usepackage{bm}
\usepackage[table]{xcolor}
\usepackage{geometry}
\geometry{
    left=2.5cm,    % 左边距
    right=2.5cm,   % 右边距
    top=2cm,       % 上边距
    bottom=2cm     % 下边距
}
% Define custom column types
\newcolumntype{L}{l}

% Revision highlighting command (matching manuscript)
\newcommand{\rev}[1]{{\color{blue}#1}}

\title{All Tables from Meta-RL LF-PID Control Revised Manuscript}
\author{Extracted from Revised Manuscript (ROB-2026-0021)}
\date{\today}

\begin{document}

\maketitle

\section*{Summary}

This document contains all tables extracted from the revised manuscript. Tables are organized by their appearance in the paper:

\begin{itemize}
    \item \textbf{Main Text Tables (10):} Core results and comparisons
    \begin{itemize}
        \item Table 1: Sensitivity of Controllability Threshold
        \item Table 2: Pilot Comparison of Optimization Algorithms
        \item Table 3: Baseline Design Summary
        \item Table 4: Performance on Franka Panda (9-DOF)
        \item Table 5: Performance on Laikago (12-DOF)
        \item Table 6: Per-Joint Tracking Error Comparison Across Platforms
        \item Table 7: Leave-One-Platform-Out Analysis
        \item Table 8: Robustness Analysis
        \item Table 9: Ablation --- Data Augmentation Impact
        \item Table 10: Optimization Ceiling Effect Quantification
    \end{itemize}
    \item \textbf{Appendix Tables (9):} Hyperparameters, configuration, and extended analysis
    \begin{itemize}
        \item Table 11: Meta-Learning Network Hyperparameters (Appendix A.1)
        \item Table 12: PPO Algorithm Hyperparameters (Appendix A.2)
        \item Table 13: Data Augmentation and LF-PID Optimization (Appendix A.3)
        \item Table 14: Reward Function and Environment (Appendix A.4)
        \item Table 15: Computing Infrastructure (Appendix A.5)
        \item Table 16: Multi-Seed Evaluation Configuration (Appendix A.6)
        \item Table 17: Computational Resource Breakdown (Appendix C)
        \item Table 18: Ablation --- RL Adaptation Subset (Appendix D)
        \item Table 19: Meta-Learning Train/Validation Performance (Appendix E)
    \end{itemize}
    \item \textbf{Nomenclature Tables (2):} Acronyms and mathematical symbols
\end{itemize}

\textbf{Total Tables: 21}

\textbf{Note:} All tables include ``Source: Authors own work'' attribution as per journal submission requirements.

% ===========================================================================
% MAIN TEXT TABLES
% ===========================================================================

\section*{MAIN TEXT TABLES}

% ===========================================================================
% TABLE 1: Sensitivity of Controllability Threshold
% ===========================================================================

\section*{Table 1: Sensitivity of Controllability Threshold}

\begin{tabular}{lcccc}
\hline
Threshold & Retained & Mean Error ($^\circ$) & NMAE (\%) & Downstream MAE ($^\circ$) \\
\hline
10$^\circ$ & 98 & 6.2$\pm$2.1 & 52.8 & 9.14 \\
20$^\circ$ & 178 & 9.8$\pm$4.7 & 48.3 & 8.02 \\
\textbf{30$^\circ$} & \textbf{232} & \textbf{13.9$\pm$8.4} & \textbf{47.1} & \textbf{7.51} \\
40$^\circ$ & 268 & 16.5$\pm$10.2 & 49.6 & 7.89 \\
No filter & 303 & 19.1$\pm$12.8 & 55.4 & 8.73 \\
\hline
\end{tabular}

\vspace{0.2cm}
\noindent
\textbf{Caption:} Sensitivity of controllability threshold on data quality and downstream meta-learning performance (Franka 7-DOF). Retained = number of virtual samples passing the filter out of 303 total. Downstream MAE is the meta-learning prediction error (degrees) on the held-out test set.

\vspace{0.2cm}
\noindent
\textbf{Source:} Authors own work.

% ===========================================================================
% TABLE 2: Pilot Comparison of Optimization Algorithms
% ===========================================================================

\section*{Table 2: Pilot Comparison of Optimization Algorithms}

\begin{tabular}{@{\extracolsep{\fill}}lcccl@{}}
\hline
\textbf{Method} & \textbf{Obj. ($^\circ$)} & \textbf{Evals} & \textbf{Time (s)} & \textbf{Notes} \\
\hline
DE+NM (ours) & 13.9$\pm$4.2 & 128 & 42$\pm$12 & Robust; no gradient needed \\
CMA-ES & 13.2$\pm$3.8 & 210 & 68$\pm$18 & Better mean; higher cost \\
L-BFGS-B & 18.7$\pm$9.1 & 85 & 28$\pm$8 & Fast; frequent local minima \\
PSO ($N$=20) & 14.5$\pm$5.0 & 400 & 125$\pm$30 & Slowest; marginal gain \\
\hline
\end{tabular}

\vspace{0.2cm}
\noindent
\textbf{Caption:} Pilot comparison of optimization algorithms on 20 representative virtual robots. All methods use identical simulation budget (2000-step trajectory evaluation). ``Obj.'' denotes the final objective $\mathcal{L}_v(\theta^*)$ in degrees. ``Evals'' denotes the number of objective function evaluations (each requiring one forward simulation).

\vspace{0.2cm}
\noindent
\textbf{Source:} Authors own work.

% ===========================================================================
% TABLE 3: Baseline Design Summary
% ===========================================================================

\section*{Table 3: Baseline Design Summary}

\begin{tabular}{@{\extracolsep{\fill}}lll}
\hline
\textbf{Baseline} & \textbf{Deploy Time} & \textbf{Key Characteristic} \\
\hline
Fuzzy-PID (manual) & hours & Interpretable scheduling, platform-dependent tuning \\
Meta-LF-PID & 0.8 ms & Instant zero-shot, no adaptation \\
Meta-LF-PID+RL & 10 min & Online refinement, additional training overhead \\
Heuristic (Z-N) & 40-120 hrs & Expert dependent, no transfer \\
Optimized (DE) & 30-60 min & High accuracy, no generalization \\
\hline
\end{tabular}

\vspace{0.2cm}
\noindent
\textbf{Source:} Authors own work.

% ===========================================================================
% TABLE 4: Performance on Franka Panda (9-DOF)
% ===========================================================================

\section*{Table 4: Performance on Franka Panda (9-DOF)}

\begin{tabular}{@{\extracolsep{\fill}}llll}
\hline
\textbf{Metric} & \textbf{Meta-LF-PID} & \textbf{Meta-LF-PID+RL} & \textbf{Improv.} \\
\hline
MAE ($^\circ$) & 7.51 & \textbf{6.26} & +16.6\% \\
RMSE ($^\circ$) & 29.32 & \textbf{25.45} & +13.2\% \\
Max Error ($^\circ$) & 48.49 & \textbf{42.12} & +13.1\% \\
Std Dev ($^\circ$) & 4.94 & \textbf{4.40} & +10.9\% \\
\hline
\end{tabular}

\vspace{0.2cm}
\noindent
\textbf{Source:} Authors own work.

% ===========================================================================
% TABLE 5: Performance on Laikago (12-DOF)
% ===========================================================================

\section*{Table 5: Performance on Laikago (12-DOF)}

\begin{tabular}{@{\extracolsep{\fill}}llll}
\hline
\textbf{Metric} & \textbf{Meta-LF-PID} & \textbf{Meta-LF-PID+RL} & \textbf{Improv.} \\
\hline
MAE ($^\circ$) & 5.91 & \textbf{5.79} & +2.1\% \\
RMSE ($^\circ$) & 29.70 & \textbf{29.29} & +1.4\% \\
Max Error ($^\circ$) & 53.09 & \textbf{50.44} & +5.0\% \\
Std Dev ($^\circ$) & 5.25 & \textbf{5.18} & +1.3\% \\
\hline
\end{tabular}

\vspace{0.2cm}
\noindent
\textbf{Source:} Authors own work.

% ===========================================================================
% TABLE 6: Per-Joint Tracking Error Comparison Across Platforms
% ===========================================================================

\section*{Table 6: Per-Joint Tracking Error Comparison Across Platforms}

\begin{tabular}{@{\extracolsep{\fill}}lllll}
\hline
\textbf{Robot} & \textbf{Joint} & \textbf{Pure Meta-LF-PID ($^\circ$)} & \textbf{Meta-LF-PID+RL ($^\circ$)} & \textbf{Improv.} \\
\hline
Franka Panda    & J1     &   2.57 &   2.26 & +12.2\% \\
                & J2     &  12.36 &   2.42 & \textbf{+80.4\%} \\
                & J3     &   4.10 &   3.87 &  +5.7\% \\
                & J4     &   6.78 &   6.49 &  +4.3\% \\
                & J5     &   5.41 &   5.32 &  +1.6\% \\
                & J6     &   4.31 &   4.19 &  +2.8\% \\
                & J7     &  11.45 &  11.26 &  +1.6\% \\
                & J8     &  10.23 &  10.19 &  +0.3\% \\
                & J9     &  10.36 &  10.33 &  +0.3\% \\
\hline
\textit{Franka Panda Avg} & &   7.51 &   6.26 & +16.6\% \\
\hline
Laikago         & J1     &   1.38 &   1.52 &  -9.6\% \\
                & J2     &   6.16 &   5.96 &  +3.3\% \\
                & J3     &  10.45 &  10.34 &  +1.0\% \\
                & J4     &   1.36 &   1.45 &  -6.8\% \\
                & J5     &   5.70 &   5.96 &  -4.5\% \\
                & J6     &  10.54 &  10.33 &  +2.0\% \\
                & J7     &   1.36 &   1.41 &  -3.1\% \\
                & J8     &   5.35 &   5.86 &  -9.6\% \\
                & J9     &  10.44 &  10.36 &  +0.8\% \\
                & J10    &   1.35 &   1.44 &  -6.7\% \\
                & J11    &   6.41 &   5.92 &  +7.7\% \\
                & J12    &  10.44 &  10.39 &  +0.5\% \\
\hline
\textit{Laikago Avg} & &   5.91 &   5.79 &  +2.1\% \\
\hline
\end{tabular}

\vspace{0.2cm}
\noindent
\textbf{Source:} Authors own work.

% ===========================================================================
% TABLE 7: Leave-One-Platform-Out Analysis
% ===========================================================================

\section*{Table 7: Leave-One-Platform-Out Analysis}

\begin{tabular}{@{\extracolsep{\fill}}lllll}
\hline
\textbf{Platform} & \textbf{Augmented} & \textbf{Training} & \textbf{Meta-LF-PID} & \textbf{Meta+RL} \\
& \textbf{Variants} & \textbf{Contrib.} & \textbf{MAE ($^\circ$)} & \textbf{MAE ($^\circ$)} \\
\hline
Franka Panda (9-DOF) & 150 & $\sim$65\% & 7.51 & 6.26 \\
KUKA iiwa (7-DOF) & 150 & $\sim$35\% & \multicolumn{2}{c}{\textit{training source only}} \\
Laikago (12-DOF) & 0 & $\leq$0.4\% & 5.91 & 5.79 \\
\hline
\end{tabular}

\vspace{0.2cm}
\noindent
\textbf{Caption:} Leave-one-platform-out analysis: training data composition vs.\ evaluation performance. Laikago contributes at most 1 base configuration sample (0.4\% of training data), making the current evaluation a \textit{de facto} leave-one-platform-out experiment for the most challenging inter-morphology case.

\vspace{0.2cm}
\noindent
\textbf{Source:} Authors own work.

% ===========================================================================
% TABLE 8: Robustness Analysis (Franka Panda)
% ===========================================================================

\section*{Table 8: Robustness Analysis (Franka Panda, MAE in $^\circ$, Representative Seed 51, 20 Episodes)}

\begin{tabular}{@{\extracolsep{\fill}}llll}
\hline
\textbf{Disturbance} & \textbf{Meta-LF-PID} & \textbf{Meta-LF-PID+RL} & \textbf{Improv.} \\
\hline
No Disturbance & 7.51 & \textbf{6.26} & +16.6\% \\
Random Force & 25.77 & \textbf{25.01} & +2.9\% \\
Payload Var. & 67.12 & \textbf{61.68} & +8.1\% \\
\textbf{Param. Uncert.} & 35.90 & \textbf{29.01} & \textbf{+19.2\%} \\
Mixed Dist. & 88.00 & \textbf{82.37} & +6.4\% \\
\hline
\textit{Average} & \textit{49.09} & \textit{44.59} & \textit{+10.0\%} \\
\hline
\end{tabular}

\vspace{0.2cm}
\noindent
\textbf{Source:} Authors own work.

% ===========================================================================
% TABLE 9: Ablation - Data Augmentation Impact
% ===========================================================================

\section*{Table 9: Ablation --- Data Augmentation Impact}

\begin{tabular}{@{\extracolsep{\fill}}lll}
\hline
\textbf{Training Data} & \textbf{Samples} & \textbf{Prediction Error (\%)} \\
\hline
Base robots only & 3 & 31.2 \\
+ Augmentation & 303 & 3.33 \\
\hline
\end{tabular}

\vspace{0.2cm}
\noindent
\textbf{Source:} Authors own work.

% ===========================================================================
% TABLE 10: Optimization Ceiling Effect Quantification
% ===========================================================================

\section*{Table 10: Optimization Ceiling Effect Quantification}

\begin{tabular}{@{\extracolsep{\fill}}lll}
\hline
\textbf{Metric} & \textbf{Franka Panda (9-DOF)} & \textbf{Laikago (12-DOF)} \\
\hline
Baseline MAE ($^\circ$)         & 7.51             & 5.91             \\
RL-adapted MAE ($^\circ$)       & 6.26             & 5.79             \\
Platform $G_{\mathrm{RL}}$ (\%)              & \textbf{16.6}    & 2.1              \\
Baseline $\mathrm{CV}_e$                     & 0.457            & 0.630            \\
Max $G_{\mathrm{RL},j}$ (\%)           & 80.4 (J2)        & 7.6 (J11)        \\
Median $G_{\mathrm{RL},j}$ (\%)         & 1.7              & $-$0.6           \\
Joints with $G_{\mathrm{RL},j} < 0$     & 0 / 9            & 6 / 12           \\
Error range ($^\circ$)          & 2.57--12.36      & 1.35--10.54      \\
\hline
\end{tabular}

\vspace{0.2cm}
\noindent
\textbf{Caption:} Quantitative characterization of the optimization ceiling effect. $G_{\mathrm{RL}}$ = platform-level RL gain rate; $\mathrm{CV}_e$ = coefficient of variation of baseline joint errors; Max $G_{\mathrm{RL},j}$ = highest per-joint gain; Joints with $G_{\mathrm{RL},j}<0$ = joints where RL slightly degrades performance.

\vspace{0.2cm}
\noindent
\textbf{Source:} Authors own work.

% ===========================================================================
% APPENDIX TABLES
% ===========================================================================

\newpage
\section*{APPENDIX TABLES}

% ===========================================================================
% TABLE 11: Meta-Learning Network Hyperparameters (Appendix A.1)
% ===========================================================================

\section*{Table 11: Meta-Learning Network Hyperparameters}

\centering
\small
\begin{tabular}{@{}p{0.42\linewidth}p{0.48\linewidth}@{}}
\hline
\textbf{Parameter} & \textbf{Value} \\
\hline
\multicolumn{2}{@{}l}{\textit{Network Architecture}} \\
Input dimension & 10 \\
Encoder layers & 256 + 256 + LayerNorm \\
Hidden layer & 128 + LayerNorm \\
Output dimension & 7 \\
Output activation & Sigmoid \\
\hline
\multicolumn{2}{@{}l}{\textit{Training Configuration}} \\
Optimizer & Adam \\
Learning rate & 0.001 \\
Weight decay & $10^{-5}$ \\
Batch size & 32 \\
Max epochs & 500 \\
Early stopping & 50 epochs \\
Loss function & Weighted MSE \\
\hline
\multicolumn{2}{@{}l}{\textit{Data Split}} \\
Training samples & 185 (80\%) \\
Validation samples & 47 (20\%) \\
Total samples & 232 \\
\hline
\multicolumn{2}{@{}l}{\textit{Training Time}} \\
Time per epoch & $\sim$1 second \\
Total time & $\sim$8 minutes \\
\hline
\end{tabular}

\vspace{0.3cm}
\noindent
\textbf{Source:} Authors own work.

\raggedright

% ===========================================================================
% TABLE 12: PPO Algorithm Hyperparameters (Appendix A.2)
% ===========================================================================

\section*{Table 12: PPO Algorithm Hyperparameters}

\centering
\small
\begin{tabular}{@{}p{0.42\linewidth}p{0.48\linewidth}@{}}
\hline
\textbf{Parameter} & \textbf{Value} \\
\hline
\multicolumn{2}{@{}l}{\textit{Network Architecture}} \\
Policy network & [s\_dim, 256, 256, 3n] \\
Value network & [s\_dim, 256, 256, 1] \\
State dimension & varies by platform \\
Action dimension & 3n (27 for Franka, 36 for Laikago) \\
Action range & $[-0.2, 0.2]$ per adjustment \\
\hline
\multicolumn{2}{@{}l}{\textit{PPO Algorithm}} \\
Total timesteps & 1,000,000 \\
Parallel envs & 8 \\
Steps per env & 2,048 \\
Batch size & 256 \\
Mini-batch size & 256 \\
Epochs & 10 \\
Clip range $\epsilon$ & 0.2 \\
\hline
\multicolumn{2}{@{}l}{\textit{Learning Rates}} \\
Policy LR & $1 \times 10^{-4}$ \\
Value LR & $1 \times 10^{-4}$ \\
LR schedule & Constant \\
\hline
\multicolumn{2}{@{}l}{\textit{GAE \& Discount}} \\
Discount $\gamma$ & 0.99 \\
GAE $\lambda$ & 0.95 \\
\hline
\multicolumn{2}{@{}l}{\textit{Loss Coefficients}} \\
Value loss coef & 0.5 \\
Entropy coef & 0.02 \\
Max grad norm & 0.5 \\
\hline
\multicolumn{2}{@{}l}{\textit{Training Time}} \\
Wall-clock time & $\sim$10 minutes \\
FPS & $\sim$1,300 \\
\hline
\end{tabular}

\vspace{0.3cm}
\noindent
\textbf{Source:} Authors own work.

\raggedright

% ===========================================================================
% TABLE 13: Data Augmentation and LF-PID Optimization (Appendix A.3)
% ===========================================================================

\section*{Table 13: Data Augmentation and LF-PID Optimization}

\centering
\small
\begin{tabular}{@{}p{0.45\linewidth}p{0.45\linewidth}@{}}
\hline
\textbf{Parameter} & \textbf{Value} \\
\hline
\multicolumn{2}{@{}l}{\textit{Property Perturbation}} \\
Mass range & $\pm 10\%$ \\
Inertia range & \rev{$\pm 15\%$} \\
Link length range & $\pm 5\%$ \\
\rev{Friction range} & \rev{$\pm 20\%$} \\
\rev{Damping range} & \rev{$\pm 30\%$} \\
Payload range & $[0, 2\times$ base$]$ \\
Virtual per robot & 100 \\
Total virtual & 300 \\
Total real & 3 \\
Before filtering & 303 \\
\hline
\multicolumn{2}{@{}l}{\textit{LF-PID Optimization}} \\
DE population & 8 \\
DE iterations & 15 \\
DE mutation $F$ & \rev{0.8} \\
DE crossover & 0.7 \\
Bounds & $K_p, K_d \in [0.1, 500]$ \\
 & $K_i \in [0, 1]$ \\
NM tolerance & $10^{-6}$ \\
Trajectory & 2000 steps (20s) \\
Time/sample & \rev{30--60s (1 core)} \\
\rev{Parallelism} & \rev{23 cores} \\
Total time & $\sim$5 min \\
\hline
\multicolumn{2}{@{}l}{\textit{Data Filtering}} \\
Error threshold & $30^\circ$ \\
Min per robot & 30 \\
Removed & 71 (23.4\%) \\
Final samples & 232 \\
\hline
\end{tabular}

\vspace{0.3cm}
\noindent
\textbf{Source:} Authors own work.

\raggedright

% ===========================================================================
% TABLE 14: Reward Function and Environment (Appendix A.4)
% ===========================================================================

\section*{Table 14: Reward Function and Environment}

\centering
\small
\begin{tabular}{@{}p{0.38\linewidth}p{0.38\linewidth}p{0.14\linewidth}@{}}
\hline
\textbf{Component} & \textbf{Formula} & \textbf{Weight} \\
\hline
Position error & $-\|q_t - q_{ref,t}\|_2$ & 1.0 \\
Velocity error & $-\|\dot{q}_t - \dot{q}_{ref,t}\|_2$ & 0.5 \\
Jerk penalty & $-\|\ddot{q}_t - \ddot{q}_{t-1}\|_2$ & 0.1 \\
LF-PID parameter change & $-\|\bm{\theta}_t - \bm{\theta}_{t-1}\|_2$ & 0.05 \\
Success bonus & +10 if $\|e_t\| < 5^\circ$ & - \\
Failure penalty & -100 if unstable & - \\
\hline
\multicolumn{3}{@{}l}{\textit{Environment Settings}} \\
\multicolumn{2}{@{}l}{Episode length} & 2000 steps \\
\multicolumn{2}{@{}l}{Control frequency} & 100 Hz \\
\multicolumn{2}{@{}l}{Simulator} & PyBullet 3.2.5 \\
\multicolumn{2}{@{}l}{Physics timestep} & 0.01 s \\
\hline
\end{tabular}

\vspace{0.3cm}
\noindent
\textbf{Source:} Authors own work.

\raggedright

% ===========================================================================
% TABLE 15: Computing Infrastructure (Appendix A.5)
% ===========================================================================

\section*{Table 15: Computing Infrastructure}

\centering
\small
\begin{tabular}{@{}p{0.4\linewidth}p{0.5\linewidth}@{}}
\hline
\textbf{Component} & \textbf{Specification} \\
\hline
\multicolumn{2}{@{}l}{\textit{Hardware}} \\
CPU & Intel i7-11800H (8c/16t) \\
GPU & NVIDIA RTX 3060 (6GB) \\
RAM & 16GB DDR4-3200 \\
OS & Ubuntu 20.04 LTS \\
\hline
\multicolumn{2}{@{}l}{\textit{Software}} \\
Python & 3.10.13 \\
PyTorch & 2.0.1 (CUDA 11.7) \\
Stable-Baselines3 & 2.0.0 \\
PyBullet & 3.2.5 \\
\hline
\multicolumn{2}{@{}l}{\textit{Training Time}} \\
Data augmentation & $\sim$5 min \\
LF-PID optimization & $\sim$5 min \\
Meta-learning & $\sim$8 min \\
RL (per robot) & $\sim$2.5 hours \\
\textbf{Total pipeline} & \textbf{$\sim$3 hours} \\
\hline
\end{tabular}

\vspace{0.3cm}
\noindent
\textbf{Source:} Authors own work.

\raggedright

% ===========================================================================
% TABLE 16: Multi-Seed Evaluation Configuration (Appendix A.6)
% ===========================================================================

\section*{Table 16: Multi-Seed Evaluation Configuration}

\centering
\small
\begin{tabular}{@{}p{0.42\linewidth}p{0.48\linewidth}@{}}
\hline
\textbf{Parameter} & \textbf{Value} \\
\hline
Seed range & 0-99 (100 different seeds) \\
Episodes per seed & 20 per disturbance scenario \\
Total episodes & 10,000 (100$\times$5$\times$20) \\
Representative seed & 51 (near-median performance) \\
Statistical metric & Mean$\pm$std across all seeds \\
\hline
\multicolumn{2}{@{}l}{\textit{Stochastic Factors Controlled by Seeds}} \\
\multicolumn{2}{@{}l}{- Trajectory initialization randomness} \\
\multicolumn{2}{@{}l}{- Disturbance timing and sequencing} \\
\multicolumn{2}{@{}l}{- Environment noise and variations} \\
\multicolumn{2}{@{}l}{- RL policy exploration randomness} \\
\hline
\end{tabular}

\vspace{0.3cm}
\noindent
\textbf{Source:} Authors own work.

\raggedright

% ===========================================================================
% TABLE 17: Computational Resource Breakdown (Appendix C)
% ===========================================================================

\section*{Table 17: Computational Resource Breakdown for All Methods}

\centering
\scriptsize
\begin{tabular}{@{}lcccl@{}}
\hline
\textbf{Method} & \textbf{Steps} & \textbf{Evals} & \textbf{Cores} & \textbf{Time} \\
\hline
\multicolumn{5}{@{}l}{\textit{Training/Optimization Phase}} \\
\quad Meta-LF-PID (offline) & 464M$^{\dagger}$ & - & 1 & 5 min \\
\quad \quad Per-sample DE+NM & 2M & 128 & 1 & 30-60s \\
\quad \quad Meta net training & - & 232$\times$500 & 1 & 3 min \\
\quad RL training (per plat.) & 8M$^{\ddagger}$ & - & 8 & 10 min \\
\quad DE optim. (per plat.) & 2M & 128 & 1 & 30-60m \\
\quad Manual Fuzzy-PID & - & - & - & 40-120h \\
\hline
\multicolumn{5}{@{}l}{\textit{Deployment/Inference Phase}} \\
\quad Meta-LF-PID infer. & - & 1 & 1 & 0.8 ms \\
\quad Meta-LF-PID+RL infer. & - & 1 & 1 & 4 ms \\
\quad DE-optimized ctrl. & - & 1 & 1 & 0.1 ms \\
\hline
\end{tabular}

\vspace{0.2cm}

\noindent \footnotesize{$^{\dagger}$Total across 232 samples during offline meta-training: $232 \times 2M = 464M$ steps.}

\noindent \footnotesize{$^{\ddagger}$RL training: $1M \times 8$ parallel environments = 8M total environment steps.}

\vspace{0.2cm}
\noindent
\normalsize\textbf{Source:} Authors own work.

\raggedright

% ===========================================================================
% TABLE 18: Ablation - RL Adaptation Subset (Appendix D)
% ===========================================================================

\section*{Table 18: Ablation --- RL Adaptation Subset Performance on Franka Panda}

\centering
\footnotesize
\begin{tabular}{@{}lcccc@{}}
\hline
\textbf{Adaptation Subset} & \textbf{Dim} & \textbf{MAE ($^\circ$)} & \textbf{Improv.} & \textbf{Stability$^*$} \\
\hline
Meta-LF-PID (no RL) & - & 7.51 & - & N/A \\
\hline
Scales only (Proposed) & 27 & \textbf{6.26} & +16.6\% & 0/5 div. \\
Consequents subset & 81 & 6.89 & +8.3\% & 2/5 div. \\
Scales + Consequents & 108 & 6.42 & +14.5\% & 3/5 div. \\
\hline
\end{tabular}

\vspace{0.2cm}

\noindent \footnotesize{$^*$Divergence across 5 seeds (42, 101, 202, 303, 404); div. = episode reward $< -200$ for $>$20\% of final 100k steps.}

\vspace{0.2cm}
\noindent
\normalsize\textbf{Source:} Authors own work.

\raggedright

% ===========================================================================
% TABLE 19: Meta-Learning Train/Validation Performance (Appendix E)
% ===========================================================================

\section*{Table 19: Meta-Learning Train/Validation Performance}

\centering
\begin{tabular}{@{}lccc@{}}
\hline
\textbf{Robot Type} & \textbf{Train (\%)} & \textbf{Val (\%)} & \textbf{Gap} \\
\hline
Franka variants (n=9) & 2.87 & 3.21 & +0.34 \\
KUKA variants (n=7) & 2.95 & 3.41 & +0.46 \\
Laikago variants (n=12) & 3.12 & 3.58 & +0.46 \\
\hline
Overall & 2.98 & 3.40 & +0.42 \\
\hline
\end{tabular}

\vspace{0.2cm}
\noindent
\normalsize\textbf{Source:} Authors own work.

\raggedright

% ===========================================================================
% NOMENCLATURE TABLES
% ===========================================================================

\newpage
\section*{NOMENCLATURE TABLES}

% ===========================================================================
% TABLE 20: Nomenclature - Acronyms
% ===========================================================================

\section*{Table 20: Nomenclature --- Acronyms}

\begin{tabular}{@{}p{0.25\linewidth}p{0.65\linewidth}@{}}
\hline
\textbf{Acronym} & \textbf{Description} \\
\hline
PID & Proportional-Integral-Derivative \\
RL & Reinforcement Learning \\
PPO & Proximal Policy Optimization \\
DE & Differential Evolution \\
MAML & Model-Agnostic Meta-Learning \\
MAE & Mean Absolute Error \\
RMSE & Root Mean Square Error \\
NMAE & Normalized Mean Absolute Error \\
DOF & Degrees of Freedom \\
MLP & Multi-Layer Perceptron \\
\hline
\end{tabular}

\vspace{0.3cm}
\noindent
\textbf{Source:} Authors own work.

% ===========================================================================
% TABLE 21: Nomenclature - Mathematical Symbols
% ===========================================================================

\section*{Table 21: Nomenclature --- Mathematical Symbols}

\begin{tabular}{@{}p{0.20\linewidth}p{0.50\linewidth}p{0.20\linewidth}@{}}
\hline
\textbf{Symbol} & \textbf{Description} & \textbf{Dim/Value} \\
\hline
\multicolumn{3}{@{}l}{\textit{Robot Features \& Network Architecture}} \\
$\mathbf{f}$ & Robot feature vector & $\mathbb{R}^{10}$ \\
$n_{dof}$ & Number of degrees of freedom & - \\
$\mathbf{h}_1, \mathbf{h}_2$ & Encoder layer outputs & $\mathbb{R}^{256}$ \\
$\mathbf{h}_{hidden}$ & Hidden layer output & $\mathbb{R}^{128}$ \\
$W_1, W_2, W_3$ & Weight matrices & - \\
$\sigma$ & Sigmoid activation & - \\[0.2cm]
\multicolumn{3}{@{}l}{\textit{LF-PID Parameters}} \\
$\bar{K}_p, \bar{K}_i, \bar{K}_d$ & Base PID gains within LF-PID & - \\
$\bm{\theta}$ & LF-PID parameter vector (base gains + scales + TS consequents) & $\mathbb{R}^{D}$ \\
$\bm{\theta}_v^*$ & Ground-truth optimal LF-PID parameters & - \\
$\hat{\bm{\theta}}_v$ & Predicted LF-PID parameters & - \\[0.2cm]
\multicolumn{3}{@{}l}{\textit{Control \& Trajectory}} \\
$q(t)$ & Joint position vector & rad \\
$\dot{q}(t)$ & Joint velocity vector & rad/s \\
$q_{ref}(t)$ & Reference trajectory & rad \\
$e(t)$ & Tracking error & rad \\
$u(t)$ & Control torque & N$\cdot$m \\
$n$ & Number of joints & - \\[0.2cm]
\multicolumn{3}{@{}l}{\textit{Loss Functions \& Optimization}} \\
$\mathcal{L}_{meta}$ & Meta-learning loss & - \\
$\mathcal{L}_v(\theta)$ & Trajectory tracking error & rad \\
$w_v$ & Sample weight & - \\
\rev{$N$} & \rev{Number of training samples (filtered)} & \rev{232} \\
$\theta^*_{global}$ & Global-search LF-PID optimum from DE & - \\
$\theta^*_v$ & Final optimized LF-PID parameters & - \\
$T$ & Trajectory length & 2000 \\[0.2cm]
\multicolumn{3}{@{}l}{\textit{RL Components}} \\
$\mathbf{s}_t$ & State observation & - \\
$\mathbf{a}_t$ & RL action (scaling factor adjustments) & $[-0.2,0.2]^{3n}$ \\
$\bm{s}$ & Per-joint input scaling factors & $\mathbb{R}^{3n}$ \\
$r_t$ & Reward signal & $[-100,10]$ \\
$\gamma$ & Discount factor & 0.99 \\
$\lambda$ & GAE parameter & 0.95 \\[0.2cm]
\multicolumn{3}{@{}l}{\textit{\color{blue}Ceiling Effect Metrics}} \\
\rev{$G_{\mathrm{RL},j}$} & \rev{Per-joint RL gain rate} & \rev{\%} \\
\rev{$G_{\mathrm{RL}}$} & \rev{Platform-level RL gain rate} & \rev{\%} \\
\rev{$\mathrm{CV}_e$} & \rev{Coefficient of variation of baseline joint errors} & \rev{-} \\
\rev{$e_{\mathrm{base},j}$} & \rev{Meta-LF-PID baseline error for joint $j$} & \rev{$^\circ$} \\
\rev{$e_{\mathrm{RL},j}$} & \rev{Meta-LF-PID+RL error for joint $j$} & \rev{$^\circ$} \\
\hline
\end{tabular}

\vspace{0.3cm}
\noindent
\textbf{Source:} Authors own work.

% ===========================================================================
% LIST OF ALL TABLES
% ===========================================================================

\newpage
\section*{List of All Tables in This Document}

\subsection*{Main Text Tables (Tables 1--10)}
\begin{enumerate}
    \item Sensitivity of Controllability Threshold
    \item Pilot Comparison of Optimization Algorithms
    \item Baseline Design Summary
    \item Performance on Franka Panda (9-DOF)
    \item Performance on Laikago (12-DOF)
    \item Per-Joint Tracking Error Comparison Across Platforms
    \item Leave-One-Platform-Out Analysis
    \item Robustness Analysis (Franka Panda, MAE in degrees)
    \item Ablation --- Data Augmentation Impact
    \item Optimization Ceiling Effect Quantification
\end{enumerate}

\subsection*{Appendix Tables (Tables 11--19)}
\begin{enumerate}
    \setcounter{enumi}{10}
    \item Meta-Learning Network Hyperparameters (Appendix A.1)
    \item PPO Algorithm Hyperparameters (Appendix A.2)
    \item Data Augmentation and LF-PID Optimization (Appendix A.3)
    \item Reward Function and Environment (Appendix A.4)
    \item Computing Infrastructure (Appendix A.5)
    \item Multi-Seed Evaluation Configuration (Appendix A.6)
    \item Computational Resource Breakdown for All Methods (Appendix C)
    \item Ablation --- RL Adaptation Subset Performance (Appendix D)
    \item Meta-Learning Train/Validation Performance (Appendix E)
\end{enumerate}

\subsection*{Nomenclature Tables (Tables 20--21)}
\begin{enumerate}
    \setcounter{enumi}{19}
    \item Nomenclature --- Acronyms
    \item Nomenclature --- Mathematical Symbols
\end{enumerate}

\vspace{0.5cm}
\noindent
\textbf{Note:} All tables preserve original LaTeX formatting from the revised manuscript, including caption labels, mathematical notation, and source attribution. Blue-colored text (\rev{like this}) indicates revisions made in response to reviewer comments.

\end{document}
